
\chapter{Фрод}\label{ch:fraud}
\highlight{Универсального признака риска нет; атака всё чаще маскируется под
нормальное поведение}. Система анализирует устройство, историю аккаунта,
частоту попыток, маршрут доставки и~поведение после оплаты и~формирует
из~этих слабых признаков общую оценку риска. Слишком строгий фильтр повышает
цену ошибочного отказа, а~слишком мягкий пропускает больше фрода.

\section{Модель угроз}\label{sec:fraud-threats}

Карточный фрод устроен как разделённый рынок: одни участники крадут данные,
другие проверяют, какие из~них ещё работают, третьи монетизируют подготовленную
атаку. Фильтр, хорошо выявляющий карточный перебор, не~выявит злоупотребление
возвратными платежами, поэтому защита начинается с~карты атак.

Самый старый класс атак~--- скимминг. Он работает, пока в~системе сохраняется
\textit{fallback} на~магнитную полосу (подробный разбор сценария~---
раздел~\ref{sec:emv-fallback}). Накладка на~терминале или банкомате считывает
данные полосы; украденные дампы (\textit{dumps}) попадают на~закрытые площадки
и~превращаются в~клон карты. Чип EMV (Europay\slash Mastercard\slash Visa,
см.~гл.~\ref{ch:emv}) резко снизил ценность такого клона: каждая операция
требует новой криптограммы, которую нельзя скопировать. Магнитная полоса
исчезла не~везде и~не~одновременно, и~пока \textit{fallback} доступен, скимминг
сохраняет силу~\cite{emvco-book1}.

Следующий класс атак~--- фрод без~предъявления карты
(\textit{card-not-present}, CNP) и~карточный перебор (\textit{carding}).
В~электронной коммерции сценарий выглядит так: у~атакующего есть номер карты
(Primary Account Number, PAN; см.~гл.~\ref{ch:tokenization}), срок действия
и~трёхзначный проверочный код на~обороте (Card Verification Value, CVV/CVV2),
полученные из~утечки, через фишинг, вредоносное ПО или на~теневом рынке. Дальше
начинается \textit{carding}: серия мелких тестовых списаний проверяет, какие
реквизиты ещё работают, а~за~успешными тестами следует крупная
покупка~\cite{mc-fraudulent-merchants-2025}. Атака оставляет плотный
повторяющийся ритм и~хорошо заметна в~частотных проверках; такие правила
срабатывают до~первой дорогой операции~\cite{visa-auth-management-2020}.

Отдельный класс атак~--- захват учётной записи (Account Takeover, ATO).
Злоумышленник входит через уже установленное доверие: массовый перебор утёкших
пар логин/пароль (\textit{credential stuffing}), фишинг, социальная инженерия,
иногда подмена SIM-карты. Он получает доступ ко~всему, что система раньше
связывала с~легитимным клиентом: истории покупок, сохранённым картам, адресам
доставки. Отдельная транзакция выглядит почти нормальной, и~поводом для
подозрения становится резкая смена устройства, географии, канала входа или
адреса доставки~\cite{mc-fraudulent-merchants-2025}.

Дружественный фрод (\textit{friendly fraud}) устроен иначе. Никто не~крадёт
карту и~не~ломает аккаунт: держатель карты сам совершает оплату,
получает товар или услугу и~затем заявляет, что операция была
несанкционированной. Дружественный фрод встречается в~цифровых товарах,
подписках и~сервисах с~нематериальной поставкой, где труднее показать момент
потребления. Без связной истории устройства, аккаунта, использования услуги
и~доставки спор превращается в~слово против~слова. Эти следы становятся
доказательной базой для~споров о~возвратных платежах
(см.~раздел~\ref{sec:disputes-evidence})~\cite{mc-first-party-trust}.

Фрод \textit{bust-out} использует накопленное доверие: клиент или продавец
сначала ведёт себя образцово, набирает доверие, увеличивает лимиты или оборот,
а~затем резко меняет профиль поведения и~выводит деньги из~системы. Для
эмитента это похоже на~кредитный обман: несколько карт, аккуратная история,
постепенный рост лимита и~затем дефолт. В~эквайринге сценарий другой:
спокойный старт, быстрый рост оборота, всё меньше возвратов и~внезапное
исчезновение продавца. Короткий период наблюдения такие случаи не~выявляет:
их видно только там, где
у~системы есть длинная история поведения и~мониторинг после
подключения~\cite{visa-merchant-screening-service,mc-fraudulent-merchants-2025}.

Набор наблюдаемых признаков меняется от~атаки к~атаке: в~одних сценариях важнее
устройство, в~других аккаунт, история доставки и~потребления или профиль
продавца. Единый фильтр на~весь этот спектр построить нельзя~--- защита
собирается из~специализированных проверок под каждый класс атак.

\section{Пороговые проверки и~статические правила}\label{sec:fraud-rules}

Первая линия защиты работает быстрее всех, потому что ей нужно принять решение
до~того, как операция превратится в~ущерб. Эту задачу решают пороговые проверки
(\textit{velocity checks}) и~статические правила. Они отвечают на~простые
вопросы: не~слишком ли много операций приходит с~одного IP, не~появилось ли
на~одном устройстве подозрительно много карт, не~отклонилась ли сумма
от~привычной истории клиента. Правило сравнивает текущую операцию
с~заранее заданным порогом и~за~миллисекунды решает, пропустить её или
отклонить. Цена простоты предсказуема: жёсткие пороги отклоняют нормальные
операции, мягкие пропускают атаку.

\highlight{Базовые правила оценивают частоту, концентрацию и~резкое отклонение
от~профиля}. Сколько операций пришло с~одного IP за~час, сколько попыток прошло
по~одной карте за~десять минут, сколько карт появилось на~одном устройстве
за~сутки, насколько текущая сумма выбивается из~профиля аккаунта.
Такие проверки дёшевы, быстро исполняются и~легко объясняются бизнесу,
комплаенсу и~регулятору.

В~коде частотная проверка сводится к~скользящему окну: события по~ключу (карта,
IP, устройство) хранятся в~очереди, события старше окна вытесняются, превышение
порога~--- отказ. Отказанные попытки тоже учитываются~--- иначе атакующий
обходит проверку, продолжая перебор сразу после достижения порога. Источник
времени вынесен в~параметр (\texttt{now}): в~рабочей системе туда передают
\texttt{time.monotonic}, а~в~тестах~--- управляемые часы.

\needspace{0.3\textheight}
\codefile{Python}{samples/ch28-velocity/velocity.py}

Атакующий умеет дробить поток под~пороги, менять IP, растягивать серию
во~времени и~подмешивать легитимный шум. Правило не~учитывает контекст: крупная
покупка в~аэропорту перед вылетом для одного клиента норма, а~для другого~---
сильная аномалия, и~правило этого не~различает.

Отдельная категория правил отвечает за~проверки карточного перебора
(\textit{carding checks}). Злоумышленник не~начинает с~дорогой покупки, поэтому
сначала появляются серии микросписаний или авторизаций на~символические суммы
с~одного IP, устройства или подсети, иногда по~десяткам разных карт. Пока атака
проверяет, какие реквизиты ещё работают, она оставляет повторяющийся ритм,
который трудно скрыть, и~именно по~нему частотные правила её выявляют.

\section{Цифровой слепок устройства и~поведенческая аналитика}\label{sec:fraud-fingerprint}

Цифровой слепок устройства (\textit{device fingerprinting})~--- набор признаков
браузера или телефона, по~сочетанию которых система распознаёт конкретное
клиентское окружение. В~него входят заголовок \texttt{User-Agent}, размер
экрана, часовой пояс, набор шрифтов, характеристики WebGL и~особенности
рендеринга \texttt{canvas} и~\texttt{audio}. Ни~один признак по~отдельности
не~решает задачу, но их комбинация даёт устойчивый профиль. Если браузер
скрывает или специально искажает эти атрибуты, точность падает, но слепок
остаётся сильным признаком риска~\cite{eckersley-fingerprinting-2010}.

Слепок становится полезнее с~накоплением истории. Устройство, которое месяцами
используется в~нормальных покупках, набирает доверие. Новый слепок, возникший
при~дорогой покупке и~с~новым адресом доставки, усиливает подозрение. На~той~же
истории строится доказательство связи клиента с~прежними операциями для~разбора
возвратных платежей.

Поведенческая аналитика добавляет данные текущей сессии~--- скорость набора,
траекторию мыши, время заполнения формы и~паузы между действиями. У~человека,
вводящего свою карту, и~у~бота, перебирающего украденные реквизиты, статические
признаки могут совпадать, а~динамика поведения расходится; анализ строит для
них разные поведенческие профили. Данные собираются пассивно и~не~усложняют
оформление покупки, тогда как браузеры с~усиленной защитой приватности
сокращают набор сигналов, а~опытный атакующий имитирует часть поведения.
Поведенческие сигналы поэтому используют только в~связке со~слепком устройства.

На~мобильных платформах для сопоставимой полноты данных нужен отдельный
комплект разработчика (Software Development Kit, SDK); часть браузеров
и~расширений специально искажает атрибуты. Слепок устройства подделать
полностью трудно, а~подделать и~при~этом не~оставить новых аномалий ещё
труднее. Этот же набор сигналов продавец передаёт в~ACS эмитента.
EMV~3-D~Secure (см.~гл.~\ref{ch:3ds}) отдельно выделяет данные об~устройстве
(\enquote{device information}), о~держателе карты и~аккаунте (\enquote{cardholder account
information}), о~признаках риска со~стороны продавца (\enquote{merchant risk
indicator}) и~о~доставке (\enquote{shipping and delivery information}) для
оценки риска на~стороне сервера управления доступом эмитента (Access Control
Server, ACS)~\cite{emvco-3ds-tech-features}.

\section{Машинное обучение в~антифроде}\label{sec:fraud-ml}

Машинное обучение (\textit{machine learning}, ML) учитывает множество слабых
признаков: время суток, сумму, географию, расстояние от~предыдущей операции,
профиль продавца, историю устройства. Там, где по~отдельности всё выглядит
нормально, модель замечает редкую комбинацию. Модель знает лишь то, чему её
научили, и~устаревает по~мере того, как меняется поведение злоумышленников,
поэтому ML дополняет правила, а~не~заменяет их.

В~реальном потоке легитимных операций гораздо больше, чем мошеннических,
поэтому простая точность (доля верных ответов) ничего не~говорит о~качестве
модели антифрода: модель, отвечающая \enquote{легитимно} на~каждый запрос,
формально выглядит почти идеальной, но бесполезна. Каждое решение модели
попадает в~одну из~четырёх клеток матрицы несоответствий (\textit{confusion
matrix}): \textbf{TP} (\textit{true positives}~--- мошеннические операции,
которые модель верно заблокировала), \textbf{FP} (\textit{false positives}~---
честные операции, ошибочно заблокированные), \textbf{TN} (\textit{true
negatives}~--- честные операции, верно пропущенные) и~\textbf{FN}
(\textit{false negatives}~--- мошеннические операции, ошибочно пропущенные).
На~этих счётчиках строятся базовые метрики:
\[
  \mathrm{precision} = \frac{\mathrm{TP}}{\mathrm{TP} + \mathrm{FP}},
  \qquad
  \mathrm{recall} \equiv \mathrm{TPR} = \frac{\mathrm{TP}}{\mathrm{TP} + \mathrm{FN}},
  \qquad
  \mathrm{FPR} = \frac{\mathrm{FP}}{\mathrm{FP} + \mathrm{TN}}.
\]
\textit{Precision}~--- доля настоящего фрода среди заблокированных; TPR
(\textit{recall}, полнота)~--- доля пойманного фрода среди всех мошеннических
операций; FPR~--- доля ложных блокировок среди всех легитимных. Компромисс
между \textit{precision} и~\textit{recall} сворачивают в~F1:
\[
  F_1 = 2 \cdot \frac{\mathrm{precision} \cdot \mathrm{recall}}
  {\mathrm{precision} + \mathrm{recall}}.
\]
Зависимость от~порога модели показывают две кривые: PR (\textit{precision}
от~\textit{recall}) с~PR-AUC и~\textbf{ROC} (Receiver Operating
Characteristic~--- TPR от~FPR) с~ROC-AUC от~0,5 у~случайного классификатора
до~1
у~идеального~\cite{sklearn-classification-metrics,roc-pr-imbalanced-2024}.
На~рис.~\ref{fig:fraud-hybrid} две рабочие точки этой ROC соответствуют
публичным порогам Stripe Radar 65 и~75~\cite{stripe-radar-risk-settings};
сами значения TPR и~FPR в~этих точках иллюстративны на~типичной CNP-кривой,
потому что \textit{Visa Advanced Authorization}, \textit{Mastercard Decision
Intelligence}, \textit{Adyen RevenueProtect}, Forter и~Riskified конкретных
чисел публично не~раскрывают. Рабочая зона FPR~2--10\,\%~--- диапазон,
который CNP-ТСП называют по~опросу MRC~2024~\cite{mrc-fraud-report-2024};
базовая доля фрода по~картам EU/EEA в~2024~году~--- 0,033\,\%
от~оборота~\cite{ecb-eba-payment-fraud-2025}. Именно из-за такой низкой
базовой доли точность блокировок (\textit{precision}) падает даже при~высоком
TPR: при~базе $0{,}033\,\%$ даже у~модели с~$\mathrm{TPR}=0{,}9$
и~$\mathrm{FPR}=0{,}02$ среди заблокированных окажется на~два порядка больше
честных операций, чем мошеннических.

\begin{figure}[htbp]
  \centering
  \includefiguresvg[width=.78\linewidth]{assets/figures/ch28-fraud/fraud-hybrid}
  \caption{ROC-кривая скоринга с~рабочими точками для~порогов
  Stripe Radar \texttt{65} и~\texttt{75} и~зоной FPR 2--10\,\%.}\label{fig:fraud-hybrid}
\end{figure}

Для производственного антифрода хватает градиентного бустинга (\textit{gradient
  boosting}~--- ансамбль решающих деревьев, в~котором каждое следующее дерево
обучается на~ошибках предыдущего): он проще в~обучении, легче объясняется
и~хорошо работает на~табличных данных. Признаки остаются простыми: окна
по~времени, агрегаты по~продавцу, скользящие суммы, расстояние от~предыдущей
транзакции. Открытая литература показывает то~же: на~табличных наборах данных
о~фроде сильными базовыми моделями остаются деревья и~градиентный бустинг,
а~качество и~отбор признаков влияют на~итоговые \textit{precision}
и~\textit{recall} не~меньше, чем выбор
архитектуры~\cite{fraud-feature-selection-2024,fraud-lightgbm-2024}.

Конкретный вектор признаков для одной транзакции CNP: сумма транзакции;
отношение суммы к~среднему чеку держателя карты за~90~дней; количество
транзакций этой карты за~последний час и~за~сутки; расстояние в~километрах
между текущим IP и~IP предыдущей операции; время с~предыдущей операции
в~минутах; возраст аккаунта у~продавца в~днях; совпадение страны банковского
идентификационного номера эмитента (Bank Identification Number, BIN) и~страны IP; совпадение хэша устройства с~ранее
виденным; код категории торговой точки (Merchant Category Code, MCC;
см.~гл.~\ref{ch:visa}) ТСП; результат проверки трёхзначного кода с~обратной
стороны карты (CVV2); результат сверки адреса держателя с~данными эмитента
(Address Verification System, AVS)\footnote{Address Verification System~---
  сервис карточных сетей, через который эмитент сравнивает почтовый индекс и/или
  цифровую часть адреса держателя, переданные продавцом в~авторизационном
запросе, со~значениями на~своей стороне~\cite{visa-avs-2024}.}; количество
отказов по~этой карте за~24~часа; количество уникальных продавцов за~последний
час; время суток, нормализованное к~часовому поясу держателя. Набор из~15--25
таких признаков покрывает большинство известных атакующих шаблонов.

Модель возвращает скор~--- нормализованную оценку в~$[0, 1]$ или её
масштабированную копию (Stripe Radar отдаёт целое 0--99, упомянутые
на~рис.~\ref{fig:fraud-hybrid} пороги соотносятся именно с~этой шкалой).
Действие назначается правилами после скоринга: при~низких значениях транзакция
проходит без~дополнительной проверки, средние направляются на~интерактивную
проверку 3DS (\enquote{3DS-challenge}) или ручную очередь, высокие
получают автоматический отказ. Пороги калибруются под целевую долю FPR,
конкретные числа оператор подбирает эмпирически по~собственным данным;
универсальных значений для произвольного потока не~существует.

В~рабочую систему порог вводят через несколько этапов: \textit{shadow mode}
(модель работает, действия не~применяются, результат сравнивается с~текущей
действующей логикой), затем A/B-эксперимент на~доле трафика, затем полный
\textit{rollout}. Прямое изменение порога без проверки в~\textit{shadow}~---
всплеск ложных блокировок или пропусков с~момента смены порога.

Главная проблема в~эксплуатации~--- дрейф концепции (\textit{concept drift}).
Мошенники подстраиваются под~то, что модель уже выявляет, и~через несколько
недель недавний признак перестаёт работать, а~давно не~переобучавшаяся модель
деградирует. Поэтому защита от~фрода устроена как непрерывный процесс: контроль
качества на~скользящем окне, регулярное
переобучение, сравнение новой и~старой версии на~отложенных
данных~\cite{concept-drift-survey-2014}.

Подтверждённая метка \enquote{это был фрод} приходит не~с~авторизации,
а~из~цикла оспаривания (см.~раздел~\ref{sec:disputes-lifecycle}):
\textit{chargeback} эмитент инициирует за~дни, но финальный исход спора
закрепляется через 60--120~дней. До~тех пор полная разметка недоступна, и~для
свежих транзакций приходится опираться на~прокси-метки: ручные проверки
операторов, отчёты держателей карт в~банк и~сигналы от~карточной сети. Поэтому
модель переобучается чаще, чем приходят финальные метки: быстрый цикл
по~прокси-меткам (дни) выявляет свежий шаблон атаки, медленный по~окончательным
меткам \textit{chargeback} (кварталы) защищает от~систематической ошибки
разметки.

Клиентскому сервису, ручным аналитикам и~бизнесу нужна понятная причина отказа,
а~ответ \enquote{модель сказала нет} не~помогает ни~восстановить платёж,
ни~улучшить систему. В~рабочих системах используют SHAP (SHapley Additive
  exPlanations~--- метод объяснения предсказаний на~основе значений Шепли
из~теории игр) и~LIME (Local Interpretable Model-agnostic Explanations~---
локальная линейная аппроксимация модели вокруг конкретной точки), показывающие
вклад отдельных признаков в~итоговое решение~\cite{lime-2016,shap-2017}.

\section{Гибридные системы: правила и~модель}\label{sec:fraud-hybrid}

\highlight{Рабочая антифрод-система гибридна: правила быстро отсекают грубые
  атаки, модель аккуратнее работает с~пограничными случаями, а~правила после
скоринга переводят численную оценку в~конкретное действие}. Внутри одного
эквайрингового потока такая схема воспроизводится у~каждого участника
антифрод-цепочки, но у~каждого свой набор данных.

У~ТСП есть точка ввода и~собственные правила предварительного отсева
(\textit{pre-screening}): запреты по~стране доставки, по~MCC, по~BIN, лимиты
на~стоимость заказа. Сам слепок устройства, на~который опирается дальнейшая
логика, собирает не~ТСП, а~скрипт PSP, вендора антифрода (Sift, Forter,
Signifyd, Kount, Riskified) или iframe ACS эмитента через 3DS Method URL,
встроенный в~страницу оплаты. У~PSP и~эквайера~--- готовая ML-модель
и~маршрутизация решений в~потоке эквайринга. У~карточной сети~---
собственные сервисы оценки риска (\textit{Visa Advanced Authorization},
\textit{Mastercard Decision Intelligence}), видящие весь трафик через сеть.
У~эмитента~--- финальный механизм оценки риска авторизации и~RBA-логика в~ACS
(см.~раздел~\ref{sec:fraud-3ds-rba}). Политика решений у~каждого своя,
а~внутренняя механика гибридной схемы у~всех устроена одинаково.

Предварительный отсев (\textit{pre-screening}) блокирует очевидные случаи:
карточный перебор, карту из~стоп-листа, явное превышение порога частоты.
Скоринговый модуль (\textit{scoring engine}) выдаёт оценку риска для операций,
прошедших первичный фильтр. Правила после скоринга (\textit{post-scoring
rules}) переводят оценку в~действие: высокие значения ведут к~блокировке,
средние требуют дополнительной верификации или проверки через 3DS, а~низкие
проходят дальше.

Допустимая задержка для каждого этапа различается. Правила
\textit{pre-screening} отрабатывают за~единицы миллисекунд, потому что любая
задержка напрямую увеличивает время ответа на~авторизацию. Скоринг ML допускает
20--50~мс, поскольку включается только после быстрого отсева. Модуль решений,
переводящий скор в~действие, вычислений не~требует и~задержки почти
не~добавляет, ограничиваясь порогами и~набором правил \textit{post-scoring}.

Правила остаются в~системе даже после того, как модель показывает хорошее
качество. Регуляторные требования напрямую обязывают блокировать определённые
категории операций по~формализованному признаку, и~модель явного правила
не~заменяет. При~разборе возвратного платежа или ответе на~запрос регулятора
нужна объяснимая причина отказа в~формулировке вида \enquote{транзакция
превысила порог по~частоте}: такая формулировка документируется и~проверяется,
а~ссылку на~решение модели регулятор не~примет.

Исходы транзакций, ручные проверки, оспаривания и~подтверждённые случаи фрода
возвращаются и~в~обучающую выборку, и~в~набор правил. Иначе модель теряет
чувствительность к~свежим атакам и~деградирует под~\textit{concept drift}.

Выбор архитектуры антифрода зависит от~объёма и~зрелости системы; единого
порога транзакций в~месяц нет. На~малом потоке собственная модель ML
нецелесообразна: обучающая выборка мала, \textit{concept drift} незаметен,
а~стоимость содержания модели сопоставима с~потерями от~фрода. Достаточно
набора правил, предоставляемых PSP
(см.~раздел~\ref{sec:ecosystem-participants}). Готовое решение от~PSP~---
например, Stripe Radar~--- поставляется как ML-модель, конфигурируемые правила
и~пороги (по~умолчанию 65 и~75 на~шкале~0--99,
см.~рис.~\ref{fig:fraud-hybrid})~\cite{stripe-radar-risk-settings}. На~среднем
потоке имеет смысл гибридная схема, в~которой готовая модель от~PSP или вендора
дополняется собственными правилами под~специфику бизнеса. На~больших объёмах
собственная модель окупается: данных достаточно для обучения, \textit{concept
drift} заметен в~реальном времени, контроль над порогами даёт прямой финансовый
эффект. Конкретные числовые пороги перехода между этими вариантами оператор
устанавливает по~собственным метрикам потерь от~фрода и~стоимости владения
моделью.

\section{Скоринг и~3DS RBA}\label{sec:fraud-3ds-rba}

Риск-ориентированная аутентификация (Risk-Based Authentication, RBA) в~3DS v2
(см.~раздел~\ref{sec:3ds-v2}) определяет, требовать ли дополнительную проверку
или разрешить авторизацию без~подтверждения. Задача продавца~--- передать
эмитенту контекст, повышающий точность решения.

В~3DS v2 запрос аутентификации содержит набор структурированных данных
от~продавца и~3DS Server. EMVCo отдельно выделяет
\enquote{cardholder account information},
\enquote{merchant risk indicator}, \enquote{shipping and delivery information}
и~\enquote{device information}; ACS собирает из~них контекст и~применяет
собственную модель~\cite{emvco-3ds,emvco-3ds-tech-features}.

Антифрод продавца вычисляет оценку риска и~передаёт её в~3DS-запросе через
атрибуты протокола. При~высоком риске продавец просит более строгую проверку,
при~низком~--- даёт эмитенту основания обойтись без~подтверждения. ACS
сопоставляет полученные данные со~своей моделью и~принимает собственное
решение; ни~продавец, ни~эмитент единолично не~управляют исходом, но
качественный контекст заметно сокращает долю как пропущенного фрода, так
и~ненужных \textit{challenge}.

На~стороне сетей и~сетевых провайдеров работают собственные сервисы оценки
риска: у~Visa это \textit{Visa Advanced Authorization} (VAA)
и~\textit{Visa Consumer Authentication Service} (VCAS), у~Mastercard~---
\textit{Decision Intelligence}. Эмитент получает от~них дополнительную оценку,
отделяющую легитимный поток от~серии тестовых и~атакующих попыток; для продавца
это ещё один сигнал в~общей авторизационной
логике~\cite{visa-auth-management-2020,mc-decision-intelligence}.

\section{Стоимость ложноположительных срабатываний}\label{sec:fraud-false-positive}
\enlargethispage{5\baselineskip}

Настройка антифрода~--- выбор между двумя ошибками матрицы несоответствий:
ложноположительным срабатыванием
(\textit{false positive}, FP~--- числитель FPR) и~ложноотрицательным
(\textit{false negative}, FN); система, не~пропускающая ни~одной атаки,
неизбежно отклоняет слишком много нормальных операций. Доля пропущенного фрода~---
\[
  1 - \mathrm{TPR} = \frac{\mathrm{FN}}{\mathrm{TP} + \mathrm{FN}}.
\]

Пропущенный фрод заметнее ошибочной блокировки: каждая мошенническая операция
даёт прямой убыток и~часто заканчивается возвратным платежом. Цена ошибочной
блокировки менее видна: упущенная выручка и~ушедший клиент не~попадают
в~ту~же отчётность, что и~\textit{chargeback}.

По данным совместного отчёта Европейского центрального банка (ЕЦБ)
и~Европейского банковского управления (European Banking Authority, EBA)
за~2024~год, доля фрода в~общем объёме карточных операций с~картами,
выпущенными в~EU/EEA, составила 0,033\,\%, а~около 83\,\% потерь в~стоимости
пришлось на~операции, инициированные удалённо (\textit{remote}, аналог
CNP)~\cite{ecb-eba-payment-fraud-2025}. По рынку России сопоставимые
агрегированные показатели публикует Центр мониторинга и~реагирования
на~компьютерные атаки в~кредитно-финансовой сфере (ФинЦЕРТ) Банка России
в~\enquote{Обзоре операций, совершённых без добровольного согласия клиентов
финансовых организаций}, причём документ охватывает все электронные средства
платежа (ЭСП), а~не~только карты~\cite{cbr-fincert-fraud-overview}. Антифрод
и~работа со~спорами связаны мониторинговыми программами схем: пропущенный фрод
увеличивает долю возвратных платежей и~переводит ТСП под надзор VAMP и~MC~ECP
(детальные пороги, формулы и~санкции~---
в~разделе~\ref{sec:disputes-monitoring}). Универсальных публичных бенчмарков
по~\textit{false positive rate} для скоринга ML и~систем только на~правилах
нет; разница между ними измеряется на~собственном потоке оператора. Каждый
лишний процентный пункт FPR обходится в
\[
  \Delta R_{\text{день}} \approx N \cdot \mathrm{AOV} \cdot \Delta_{\mathrm{FPR}},
\]
где $N$~--- объём транзакций в~день, $\mathrm{AOV}$ (Average Order Value)~---
средний чек, $\Delta_{\mathrm{FPR}}$~--- приращение доли ложных блокировок
(для одного процентного пункта~--- $0{,}01$).

Когда легитимная транзакция ошибочно блокируется, клиент обращается
в~поддержку, а~продавец теряет текущую оплату и~будущие покупки этого человека.
Цена ошибки выше, если отказ произошёл при~первой покупке: такой клиент
не~возвращается~\cite{visa-auth-management-2020}.

В~детекции фрода точность и~полнота считаются для класса мошеннических
операций; общий поток для этих метрик не~подходит. Чем ниже порог, тем больше
блокировок и~меньше фрода, но больше ошибочных отказов; повышение порога даёт
обратный эффект. Оптимальная точка зависит от~экономики продукта. Для цифровых
товаров с~низкой себестоимостью доставки фильтр строже, а~для физической
доставки и~дорогого возврата порог приходится смягчать.

Платёжный шлюз с~высокой долей цифровой подписки поддерживает FPR
на~уровне единиц процентов или ниже, потому что цена пропущенного фрода~---
немедленный \textit{chargeback}, а~стоимость отказа~--- минута поддержки
и~новая попытка.
Маркетплейс физических товаров с~доставкой и~возвратом за~счёт продавца
выбирает FPR ближе к~верхней границе диапазона MRC~2024
(см.~рис.~\ref{fig:fraud-hybrid}): ошибочная блокировка экономит ему складские
издержки на~возврат, а~средний чек выше. Конкретные значения~--- авторская
оценка типовой настройки; публично нормированных подсегментных бенчмарков рынок
не~раскрывает. Стоимость каждой ошибки команда антифрода переводит в~рубли
и~считает на~собственном потоке: упущенная маржа $m \cdot \mathrm{AOV}$
от~ложной блокировки (где $m$~--- доля валовой маржи в~чеке) против полной
стоимости пропущенного фрода~--- возврата $\mathrm{AOV}$, комиссии процессора
за~возвратный платёж (\textit{chargeback fee}, типично
15--100~USD~\cite{chargeback-fee-2024}), потерянного товара и~доставки, плюс
вклад в~\textit{ratio} VAMP/MC~ECP, после превышения которого
платёжная схема начисляет штраф. Точка безразличия
(\textit{indifference point}~--- равенство стоимостей ложной
блокировки и~пропущенного фрода) даёт верхний предел разумного FPR, выше
которого настройка фильтра уже не~окупается даже при~нулевом пропуске атак.

\practicenote{%
  Стоимость \textit{false positive} считают явно, складывая упущенную
  транзакцию, риск потерять клиента и~расходы поддержки. Такой расчёт выявляет
  случаи, когда система настроена агрессивнее, чем бизнесу выгодно.

}
